\documentclass[letterpaper,11pt]{article}

\usepackage{latexsym}
\usepackage[empty]{fullpage}
\usepackage{titlesec}
\usepackage{marvosym}
\usepackage[usenames,dvipsnames]{color}
\usepackage{verbatim}
\usepackage{enumitem}
\usepackage[hidelinks]{hyperref}
\usepackage{fancyhdr}
\usepackage[english]{babel}
\usepackage[utf8]{inputenc}
\usepackage{textgreek}
\usepackage{tabularx}
\input{glyphtounicode}

%----------FONT OPTIONS----------
% sans-serif
% \usepackage[sfdefault]{FiraSans}
% \usepackage[sfdefault]{roboto}
% \usepackage[sfdefault]{noto-sans}
% \usepackage[default]{sourcesanspro}

% serif
% \usepackage{CormorantGaramond}
% \usepackage{charter}


\pagestyle{fancy}
\fancyhf{} % clear all header and footer fields
\fancyfoot{}
\renewcommand{\headrulewidth}{0pt}
\renewcommand{\footrulewidth}{0pt}

% Adjust margins
\addtolength{\oddsidemargin}{-0.5in}
\addtolength{\evensidemargin}{-0.5in}
\addtolength{\textwidth}{1in}
\addtolength{\topmargin}{-.5in}
\addtolength{\textheight}{1.0in}

\urlstyle{same}

\raggedbottom
\raggedright
\setlength{\tabcolsep}{0in}

% Sections formatting
\titleformat{\section}{
  \vspace{-4pt}\scshape\raggedright\large
}{}{0em}{}[\color{black}\titlerule \vspace{-5pt}]

% Ensure that generate pdf is machine readable/ATS parsable
\pdfgentounicode=1

%-------------------------
% Custom commands
\newcommand{\resumeItem}[1]{
  \item\small{
    {#1 \vspace{-2pt}}
  }
}

\newcommand{\resumeSubheading}[4]{
  \vspace{-2pt}\item
    \begin{tabular*}{0.97\textwidth}[t]{l@{\extracolsep{\fill}}r}
      \textbf{#1} & #2 \\
      \textit{\small#3} & \textit{\small #4} \\
    \end{tabular*}\vspace{-7pt}
}

\newcommand{\resumeSubSubheading}[2]{
    \item
    \begin{tabular*}{0.97\textwidth}{l@{\extracolsep{\fill}}r}
      \textit{\small#1} & \textit{\small #2} \\
    \end{tabular*}\vspace{-7pt}
}

\newcommand{\resumeProjectHeading}[2]{
    \item
    \begin{tabular*}{0.97\textwidth}{l@{\extracolsep{\fill}}r}
      \small#1 & #2 \\
    \end{tabular*}\vspace{-7pt}
}

\newcommand{\resumePatentHeading}[3]{
    \item
    \begin{tabularx}{0.97\textwidth}[t]{@{}>{\raggedright\arraybackslash}X r@{}}
      \small\textbf{#1} -- \textit{#2} & \small #3 \\
    \end{tabularx}\vspace{-7pt}
}

\newcommand{\resumeSubItem}[1]{\resumeItem{#1}\vspace{-4pt}}

\renewcommand\labelitemii{$\vcenter{\hbox{\tiny$\bullet$}}$}

\newcommand{\resumeSubHeadingListStart}{\begin{itemize}[leftmargin=0.15in, label={}]}
\newcommand{\resumeSubHeadingListEnd}{\end{itemize}}
\newcommand{\resumeItemListStart}{\begin{itemize}}
\newcommand{\resumeItemListEnd}{\end{itemize}\vspace{-5pt}}

%-------------------------------------------
%%%%%%  RESUME STARTS HERE  %%%%%%%%%%%%%%%%%%%%%%%%%%%%

%-------------------------------------------
\begin{document}
 
%----------HEADING----------
\begin{center}
    \textbf{\Huge \scshape Harinandan Kotamsetti} \\ \vspace{1pt}
    \small Cupertino, CA $|$ +1 (669) 264-9458 $|$ \href{mailto:hareee234@gmail.com}{\underline{hareee234@gmail.com}} $|$
    \href{https://www.linkedin.com/in/hareee234/}{\underline{linkedin.com/in/hareee234}} $|$
    \href{https://github.com/AwesomenessReborn}{\underline{github.com/AwesomenessReborn}}
\end{center}
 
%-----------EDUCATION-----------
\section{Education}
  \resumeSubHeadingListStart
    \resumeSubheading
      {\textbf{San Jose State University}}{San Jose, CA}
      {Bachelor of Science in Computer Engineering}{Expected Dec 2026}
      \resumeItemListStart
        \resumeItem{\textbf{Coursework:} OS, Computer Architecture, C/C++, FPGA, Verilog, Linear Algebra, x86 \& MIPS Assembly}
        
      \resumeItemListEnd
  \resumeSubHeadingListEnd
 
%-----------EXPERIENCE-----------
\section{Experience}
  \resumeSubHeadingListStart
 
    \resumeSubheading
      {\textbf{STMicroelectronics}}{Santa Clara, CA}
      {\textbf{Software \& Algorithm Intern}}{Mar 2025 -- Present}
      \resumeItemListStart
      
        \resumeItem{\textbf{Implemented a C++ JNI inference bridge} for a \textbf{129 KB, 103K-parameter CNN-LSTM} targeting Qualcomm Snapdragon NPU via \textbf{QNN/QAIRT C APIs}, connecting native model execution and preprocessing to a Kotlin/Compose Android app; \textbf{0.059 m/s MAE} on IMU walking-speed estimation.}
        
        \resumeItem{\textbf{Ported IMU preprocessing from Python to C++} over \textbf{512-sample windows at 208 Hz}, using Butterworth filtering, RobustScaler normalization, and ring-buffer windowing, validated through C/Python tests.}    
        
        \resumeItem{\textbf{Adapted an XR head-motion forecasting model for edge-accelerator constraints}, profiling parameter/context-size tradeoffs and redesigning incompatible attention operations when Qualcomm accelerator kernels could not support the original compute graph.}
        
        \resumeItem{\textbf{Deployed always-on sensor intelligence} at \textbf{50 Hz entirely in-sensor} using ISPU/MLC with zero host-MCU compute for fall detection and personalized wrist-tilt recognition on SensorTile.box.}
                
      \resumeItemListEnd
 
    \resumeSubheading
      {\textbf{Spartan Racing (Formula SAE)}}{San Jose, CA}
      {\textbf{Embedded Systems \& Software Developer}}{Aug 2023 -- Feb 2025}
      \resumeItemListStart
        \resumeItem{\textbf{Improved STM32F105 BMS round-robin polling latency (embedded C firmware)} from 7 s to 0.5 s for a 400 V, 12S8P x 8 module Formula SAE battery by removing blocking CAN/thermistor delays, refactoring module reads, and replacing dynamic buffers with static lifetime buffers.}
        
        \resumeItem{\textbf{Implemented interrupt-driven BMS fault handling} using \textbf{TIM7/NVIC} \textbf{ISR} logic and 8-bit fault/ 
        warning state bytes for HV safety shutdown signaling across an 8-module accumulator architecture.}
        
        \resumeItem{\textbf{Improved CAN telemetry reliability} by replacing blocking 10 ms transmit delays with STM32 3-mailbox availability checks and 0.5 s reconnect retry logic for reliable 1 Hz BMS data transmission.}
        
        \resumeItem{\textbf{Refactored LTC6811 BMS firmware} \textbf{for LTC6813 migration}, updating ADC register handling, static buffer sizing, and CAN data for 16-bit cell voltage/temperature telemetry across the \textbf{768-cell pack} architecture.}
        
      \resumeItemListEnd
 
  \resumeSubHeadingListEnd
 
%-----------PATENTS & INVENTION DISCLOSURES-----------
\section{Patents \& Invention Disclosures}
  \resumeSubHeadingListStart
    \resumePatentHeading
      {Glasses with Wellness and Health Monitoring}
      {STMicroelectronics, Co-inventor}
      {Filed May 2026}
    \resumePatentHeading
      {Motion-to-Photon Latency Compensation Using Time-Series Foundation Models for XR/AR/VR Glasses}
      {STMicroelectronics, Co-inventor}
      {Invention Disclosure, 2026}
  \resumeSubHeadingListEnd
 
%-----------PROJECTS-----------
\section{Projects}
    \resumeSubHeadingListStart
 
      \resumeProjectHeading
          {\textbf{MIPS SoC with Hardware Accelerator} $|$ \emph{Verilog, Xilinx Vivado, Basys 3 FPGA}}{2024}
          \resumeItemListStart
            \resumeItem{Designed and deployed \textbf{32-bit MIPS SoC} on Artix-7 FPGA with custom single-cycle CPU, memory-mapped GPIO, and hardware factorial accelerator using Moore FSM with datapath/control separation.}
            
          \resumeItemListEnd
 
      \resumeProjectHeading
          {\textbf{MERIT UAV --- Autonomous Surveillance Drone} $|$ \emph{Python, pymavlink, YOLOv8, CUDA, Jetson}}{2026}
          \resumeItemListStart
            \resumeItem{Developed \textbf{MAVLink flight controller integration} on NVIDIA Jetson Orin Nano over soldered UART to Pixhawk in GUIDED mode, with a concurrent GPS telemetry daemon and YOLOv8n CUDA inference pipeline for geotagged detection; flight control validated on hardware, vision pipeline ground-tested.}
            
          \resumeItemListEnd 
 
    \resumeSubHeadingListEnd
 
%-----------TECHNICAL SKILLS-----------
\section{Technical Skills}
 \begin{itemize}[leftmargin=0.15in, label={}]
    \small{\item{
     \textbf{Languages}{: C, C++, Python, Verilog, Java, Kotlin, x86/MIPS Assembly} \\
     \textbf{Embedded \& Protocols}{: STM32 (ARM Cortex-M3), Device Drivers (SPI, isoSPI, I2C, CAN, BLE), MEMS IMUs, Real-Time Systems, DSP, Signal Processing} \\
     \textbf{Testing \& Validation}{: HIL Testing, JNI Harness, Regression Testing, Software Simulation, Performance Benchmarking} \\
     \textbf{EDA \& Simulation}{: Altium, LTspice, MATLAB, Simulink, Xilinx Vivado} \\
     \textbf{Development Tools}{: Git, JTAG Debuggers, Oscilloscopes, CMake, Make, STM32CubeIDE} \\
     \textbf{ML/AI \& Data}{: TensorFlow, TFLite, PyTorch, ONNX, SciPy, InfluxDB, Grafana, Docker}
    }}
 \end{itemize}
 
%-------------------------------------------
\end{document}
